\documentclass[11pt]{article}
\usepackage{amsmath, amssymb, geometry, amsthm}
\geometry{a4paper, margin=1in}

\title{Problem Sets: Properties of Determinants}
\author{}
\date{}

\begin{document}

\maketitle

\tableofcontents
\newpage

\section{Concept Checklist \& Question Mapping}
This section outlines the core concepts and problem archetypes derived from the text, mapped to the corresponding questions in the problem set.

\begin{itemize}
    \item \textbf{C1: Determinant via Row Reduction (Type 1)} - Using row replacements and swaps to reach upper triangular form. (Questions 1, 2, 3)
    \item \textbf{C2: Hybrid Evaluation (Type 5)} - Combining one or two row operations with cofactor expansion to minimize arithmetic. (Questions 4, 5, 6)
    \item \textbf{C3: Abstract Row Operations (Type 2)} - Deducing the determinant of a modified matrix by tracking row operations backward. (Questions 7, 8, 9, 10)
    \item \textbf{C4: Invertibility via Determinants (Type 3)} - Using $\det(A) \neq 0$ to determine if a matrix is invertible, often solving for a parameter. (Questions 11, 12)
    \item \textbf{C5: Linear Independence via Determinants (Type 3)} - Checking if column vectors form a linearly independent set by evaluating the determinant. (Questions 13, 14)
    \item \textbf{C6: Multiplicative Property (Type 4)} - Applying $\det(AB) = \det(A)\det(B)$. (Questions 15, 21)
    \item \textbf{C7: Transpose Property (Type 4)} - Applying $\det(A^T) = \det(A)$. (Question 18)
    \item \textbf{C8: Inverse Property (Type 4)} - Applying $\det(A^{-1}) = 1/\det(A)$. (Questions 17, 18)
    \item \textbf{C9: Scalar Multiplication of a Matrix (Type 4)} - Applying $\det(cA) = c^n \det(A)$. (Questions 16, 18, 19)
    \item \textbf{C10: Determinants of Block Matrices (Type 6)} - Evaluating determinants of matrices partitioned into triangular blocks. (Questions 23, 24)
    \item \textbf{C11: The Vandermonde Determinant (Type 7)} - Recognizing and evaluating Vandermonde matrices. (Questions 25, 26)
    \item \textbf{C12: Theoretical Proofs \& Flaw Finding (Type 8)} - Identifying common traps (linearity trap, scalar trap, sign flip) and proving abstract properties (nilpotent, skew-symmetric). (Questions 20, 22, 27, 28, 29, 30)
\end{itemize}

\newpage

\section{Problem Sets}

\begin{enumerate}

    % C1: Row Reduction
    \item \textbf{[C1]} Find the determinant of the following matrix by row reduction to echelon form:
    \[ A = \begin{bmatrix} 1 & 3 & 2 \\ 2 & 7 & 5 \\ -1 & -2 & 1 \end{bmatrix} \]

    \item \textbf{[C1]} Evaluate the determinant using row operations, being careful with row swaps:
    \[ B = \begin{bmatrix} 0 & 2 & -4 \\ 3 & 1 & 5 \\ 6 & 2 & 12 \end{bmatrix} \]

    \item \textbf{[C1]} Compute the determinant of the matrix:
    \[ C = \begin{bmatrix} 4 & 1 & 7 & -2 \\ 0 & -3 & 5 & 1 \\ 0 & 0 & 2 & 8 \\ 0 & 0 & 0 & -1 \end{bmatrix} \]

    % C2: Hybrid Evaluation
    \item \textbf{[C2]} Combine the methods of row reduction and cofactor expansion to compute the determinant:
    \[ D = \begin{bmatrix} 3 & 2 & -1 & 4 \\ 2 & 0 & 1 & -2 \\ -4 & 0 & 3 & 1 \\ 5 & -2 & 2 & 3 \end{bmatrix} \]

    \item \textbf{[C2]} Use the hybrid method to evaluate:
    \[ E = \begin{bmatrix} 1 & 0 & 2 & -1 \\ 3 & 0 & 0 & 4 \\ 2 & 1 & 4 & -3 \\ 1 & 0 & 5 & 0 \end{bmatrix} \]

    \item \textbf{[C2]} Find the determinant by creating zeros in a strategic row or column first:
    \[ F = \begin{bmatrix} 2 & 1 & 3 & 1 \\ 0 & 4 & 0 & 0 \\ 1 & 2 & -1 & 2 \\ 3 & 1 & 2 & -1 \end{bmatrix} \]

    % C3: Abstract Row Operations
    \item \textbf{[C3]} Let $A$ be a $3 \times 3$ matrix with $\det(A) = 7$. Matrix $B$ is obtained from $A$ by swapping Row 1 and Row 3, and then replacing Row 2 with Row 2 minus 4 times Row 1. Find $\det(B)$.

    \item \textbf{[C3]} Given $\begin{vmatrix} a & b & c \\ d & e & f \\ g & h & i \end{vmatrix} = -3$, find the value of $\begin{vmatrix} a & b & c \\ 2d & 2e & 2f \\ g+3a & h+3b & i+3c \end{vmatrix}$.

    \item \textbf{[C3]} Given $\det(A) = 4$ for a $3 \times 3$ matrix $A = \begin{bmatrix} a & b & c \\ d & e & f \\ g & h & i \end{bmatrix}$, evaluate $\begin{vmatrix} 3a & 3b & 3c \\ d & e & f \\ -g & -h & -i \end{vmatrix}$.

    \item \textbf{[C3]} Without fully expanding, determine the value of $\begin{vmatrix} a & b & c \\ d-2a & e-2b & f-2c \\ 4a & 4b & 4c \end{vmatrix}$ and justify your answer.

    % C4: Invertibility
    \item \textbf{[C4]} For what values of $k$ is the following matrix invertible?
    \[ G = \begin{bmatrix} 1 & -1 & 2 \\ 2 & k & 4 \\ 3 & -3 & k^2 \end{bmatrix} \]

    \item \textbf{[C4]} Use determinants to decide if the matrix $H$ is invertible:
    \[ H = \begin{bmatrix} 2 & 4 & 2 \\ 1 & 2 & 1 \\ 5 & 7 & 9 \end{bmatrix} \]

    % C5: Linear Independence
    \item \textbf{[C5]} Use determinants to decide if the following set of vectors in $\mathbb{R}^3$ is linearly independent:
    \[ \mathbf{v}_1 = \begin{bmatrix} 1 \\ -1 \\ 2 \end{bmatrix}, \quad \mathbf{v}_2 = \begin{bmatrix} 2 \\ 1 \\ 3 \end{bmatrix}, \quad \mathbf{v}_3 = \begin{bmatrix} 0 \\ -3 \\ 1 \end{bmatrix} \]

    \item \textbf{[C5]} Find the value(s) of $h$ for which the vectors are linearly dependent:
    \[ \mathbf{u}_1 = \begin{bmatrix} 1 \\ 2 \\ 3 \end{bmatrix}, \quad \mathbf{u}_2 = \begin{bmatrix} 2 \\ 3 \\ 4 \end{bmatrix}, \quad \mathbf{u}_3 = \begin{bmatrix} 3 \\ 4 \\ h \end{bmatrix} \]

    % C6-C9: Algebraic Properties
    \item \textbf{[C6]} Let $A$ and $B$ be $n \times n$ matrices with $\det(A) = 4$ and $\det(B) = -3$. Find $\det(AB)$.

    \item \textbf{[C9]} Let $A$ be a $3 \times 3$ matrix with $\det(A) = 5$. Find $\det(2A)$.

    \item \textbf{[C8]} Let $A$ be an invertible $4 \times 4$ matrix with $\det(A) = -2$. Find $\det(A^{-1})$.

    \item \textbf{[C7, C8, C9]} Let $A$ and $B$ be $3 \times 3$ matrices with $\det(A) = 3$ and $\det(B) = 6$. Compute $\det(2 A^T B^{-1})$.

    \item \textbf{[C9]} Let $A$ be a $5 \times 5$ matrix with $\det(A) = 7$. Find $\det(-A)$.

    \item \textbf{[C12]} Let $A$ be an $n \times n$ matrix with real entries such that $\det(A^4) = 16$. What are the possible values for $\det(A)$?

    \item \textbf{[C6, C8]} Let $A$ and $P$ be $n \times n$ matrices, where $P$ is invertible. Prove that $\det(P^{-1}AP) = \det(A)$.

    \item \textbf{[C12]} An orthogonal matrix $Q$ is a square matrix such that $Q^T Q = I$. What are the possible values for $\det(Q)$?

    % C10: Block Matrices
    \item \textbf{[C10]} Compute the determinant of the block matrix:
    \[ M = \begin{bmatrix} 3 & 2 & 9 & 1 \\ 0 & 4 & 8 & 4 \\ 0 & 0 & 2 & 5 \\ 0 & 0 & 1 & 3 \end{bmatrix} \]

    \item \textbf{[C10]} Compute the determinant of the block matrix:
    \[ N = \begin{bmatrix} 3 & 0 & 0 \\ 5 & 1 & 2 \\ -1 & 3 & 4 \end{bmatrix} \]

    % C11: Vandermonde
    \item \textbf{[C11]} Evaluate the Vandermonde determinant:
    \[ V = \begin{vmatrix} 1 & 2 & 4 \\ 1 & 3 & 9 \\ 1 & 5 & 25 \end{vmatrix} \]

    \item \textbf{[C11]} Without expanding, explain why the determinant of the following matrix is zero:
    \[ W = \begin{bmatrix} 1 & 1 & 1 \\ 1 & 1 & 4 \\ 1 & 1 & 16 \end{bmatrix} \]

    % C12: Flaws and Proofs
    \item \textbf{[C12]} \textbf{Find the Flaw:} Let $A = \begin{bmatrix} 2 & 0 \\ 0 & 2 \end{bmatrix}$ and $B = \begin{bmatrix} 3 & 0 \\ 0 & 3 \end{bmatrix}$. 
    \textit{Flawed Solution:} By the linearity property, $\det(A+B) = \det(A) + \det(B) = 4 + 9 = 13$. Identify the error and provide the correct solution.

    \item \textbf{[C12]} A matrix $A$ is called skew-symmetric if $A^T = -A$. Prove that if $A$ is an $n \times n$ skew-symmetric matrix and $n$ is odd, then $\det(A) = 0$.

    \item \textbf{[C12]} A square matrix $A$ is nilpotent if $A^k = 0$ for some positive integer $k$. Prove that if $A$ is nilpotent, then $\det(A) = 0$.

    \item \textbf{[C12]} \textbf{Find the Flaw:} Evaluate $\det(A)$ where $A = \begin{bmatrix} 1 & 2 \\ 3 & 4 \end{bmatrix}$ using row operations.
    \textit{Flawed Solution:} Apply $R_2 \leftarrow 3R_1 - R_2$ to get $\begin{bmatrix} 1 & 2 \\ 0 & 2 \end{bmatrix}$. The determinant of this upper triangular matrix is $1 \times 2 = 2$. Identify the error and provide the correct solution.

\end{enumerate}

\newpage

\section{Detailed Solutions}

\begin{enumerate}

    % Solution 1
    \item \textbf{Solution:} We use row replacement operations to create zeros below the first pivot.
    \[ \begin{vmatrix} 1 & 3 & 2 \\ 2 & 7 & 5 \\ -1 & -2 & 1 \end{vmatrix} \]
    Apply $R_2 \leftarrow R_2 - 2R_1$ and $R_3 \leftarrow R_3 + R_1$:
    \[ = \begin{vmatrix} 1 & 3 & 2 \\ 0 & 1 & 1 \\ 0 & 1 & 3 \end{vmatrix} \]
    Now, create a zero below the second pivot using $R_3 \leftarrow R_3 - R_2$:
    \[ = \begin{vmatrix} 1 & 3 & 2 \\ 0 & 1 & 1 \\ 0 & 0 & 2 \end{vmatrix} \]
    The matrix is in upper triangular form. The determinant is the product of the diagonal entries: $(1)(1)(2) = 2$.
    \textbf{Final Answer:} $2$

    % Solution 2
    \item \textbf{Solution:} The top-left entry is 0, so we must perform a row swap. Swapping $R_1$ and $R_2$ multiplies the determinant by $-1$.
    \[ \begin{vmatrix} 0 & 2 & -4 \\ 3 & 1 & 5 \\ 6 & 2 & 12 \end{vmatrix} = - \begin{vmatrix} 3 & 1 & 5 \\ 0 & 2 & -4 \\ 6 & 2 & 12 \end{vmatrix} \]
    Create a zero in the bottom left using $R_3 \leftarrow R_3 - 2R_1$:
    \[ = - \begin{vmatrix} 3 & 1 & 5 \\ 0 & 2 & -4 \\ 0 & 0 & 2 \end{vmatrix} \]
    The matrix is now upper triangular. Multiply the diagonal entries and the $-1$ outside:
    $\det(B) = -(3 \times 2 \times 2) = -12$.
    \textbf{Final Answer:} $-12$

    % Solution 3
    \item \textbf{Solution:} The matrix $C$ is already in upper triangular form. By the Triangular Shortcut, the determinant is simply the product of its main diagonal entries.
    $\det(C) = (4)(-3)(2)(-1) = 24$.
    \textbf{Final Answer:} $24$

    % Solution 4
    \item \textbf{Solution:} We want to create zeros in a column or row to make cofactor expansion easy. Notice Column 2 already has two zeros. We can eliminate the $-2$ in Row 4 by using the $2$ in Row 1.
    Apply $R_4 \leftarrow R_4 + R_1$:
    \[ \begin{vmatrix} 3 & 2 & -1 & 4 \\ 2 & 0 & 1 & -2 \\ -4 & 0 & 3 & 1 \\ 8 & 0 & 1 & 7 \end{vmatrix} \]
    Expand along Column 2. The only non-zero term is the $2$ at position $(1,2)$. The sign is $(-1)^{1+2} = -1$.
    \[ = 2(-1) \begin{vmatrix} 2 & 1 & -2 \\ -4 & 3 & 1 \\ 8 & 1 & 7 \end{vmatrix} = -2 \begin{vmatrix} 2 & 1 & -2 \\ -4 & 3 & 1 \\ 8 & 1 & 7 \end{vmatrix} \]
    Now, use row operations on the $3 \times 3$ matrix. Let's eliminate entries in Column 2 using $R_1$.
    $R_2 \leftarrow R_2 - 3R_1$ and $R_3 \leftarrow R_3 - R_1$:
    \[ = -2 \begin{vmatrix} 2 & 1 & -2 \\ -10 & 0 & 7 \\ 6 & 0 & 9 \end{vmatrix} \]
    Expand along Column 2 of this new matrix. The sign for position $(1,2)$ is negative.
    \[ = -2 \left( 1(-1) \begin{vmatrix} -10 & 7 \\ 6 & 9 \end{vmatrix} \right) = 2 ((-10)(9) - (7)(6)) = 2(-90 - 42) = 2(-132) = -264 \]
    \textbf{Final Answer:} $-264$

    % Solution 5
    \item \textbf{Solution:} Column 2 has three zeros. We immediately expand along Column 2. The non-zero entry is $1$ at position $(3,2)$. The sign is $(-1)^{3+2} = -1$.
    \[ \det(E) = 1(-1) \begin{vmatrix} 1 & 2 & -1 \\ 3 & 0 & 4 \\ 1 & 5 & 0 \end{vmatrix} = - \begin{vmatrix} 1 & 2 & -1 \\ 3 & 0 & 4 \\ 1 & 5 & 0 \end{vmatrix} \]
    Expand the $3 \times 3$ matrix along Row 2:
    \[ = - \left( 3(-1)^{2+1} \begin{vmatrix} 2 & -1 \\ 5 & 0 \end{vmatrix} + 4(-1)^{2+3} \begin{vmatrix} 1 & 2 \\ 1 & 5 \end{vmatrix} \right) \]
    \[ = - \left( -3(0 - (-5)) - 4(5 - 2) \right) = - \left( -3(5) - 4(3) \right) = -(-15 - 12) = -(-27) = 27 \]
    \textbf{Final Answer:} $27$

    % Solution 6
    \item \textbf{Solution:} Row 2 has three zeros. Expand along Row 2. The non-zero entry is $4$ at position $(2,2)$. The sign is $(-1)^{2+2} = +1$.
    \[ \det(F) = 4 \begin{vmatrix} 2 & 3 & 1 \\ 1 & -1 & 2 \\ 3 & 2 & -1 \end{vmatrix} \]
    Evaluate the $3 \times 3$ determinant (e.g., by cofactor expansion along Row 1):
    \[ = 4 \left( 2 \begin{vmatrix} -1 & 2 \\ 2 & -1 \end{vmatrix} - 3 \begin{vmatrix} 1 & 2 \\ 3 & -1 \end{vmatrix} + 1 \begin{vmatrix} 1 & -1 \\ 3 & 2 \end{vmatrix} \right) \]
    \[ = 4 \left( 2(1 - 4) - 3(-1 - 6) + 1(2 - (-3)) \right) = 4 \left( 2(-3) - 3(-7) + 1(5) \right) \]
    \[ = 4(-6 + 21 + 5) = 4(20) = 80 \]
    \textbf{Final Answer:} $80$

    % Solution 7
    \item \textbf{Solution:} Swapping two rows multiplies the determinant by $-1$. Replacing a row with itself plus a multiple of another row does not change the determinant.
    Therefore, $\det(B) = -1 \times \det(A) = -1 \times 7 = -7$.
    \textbf{Final Answer:} $-7$

    % Solution 8
    \item \textbf{Solution:} We work backward from the modified matrix.
    First, apply $R_3 \leftarrow R_3 - 3R_1$. This replacement operation does not change the determinant:
    \[ \begin{vmatrix} a & b & c \\ 2d & 2e & 2f \\ g & h & i \end{vmatrix} \]
    Next, factor out the scalar $2$ from the second row. By the scalar property of determinants, factoring a scalar from a single row pulls it out to the front:
    \[ = 2 \begin{vmatrix} a & b & c \\ d & e & f \\ g & h & i \end{vmatrix} = 2(-3) = -6 \]
    \textbf{Final Answer:} $-6$

    % Solution 9
    \item \textbf{Solution:} Factor out $3$ from Row 1, and factor out $-1$ from Row 3:
    \[ \begin{vmatrix} 3a & 3b & 3c \\ d & e & f \\ -g & -h & -i \end{vmatrix} = 3 \begin{vmatrix} a & b & c \\ d & e & f \\ -g & -h & -i \end{vmatrix} = 3(-1) \begin{vmatrix} a & b & c \\ d & e & f \\ g & h & i \end{vmatrix} \]
    Substitute $\det(A) = 4$: $3(-1)(4) = -12$.
    \textbf{Final Answer:} $-12$

    % Solution 10
    \item \textbf{Solution:} First, apply the row replacement $R_2 \leftarrow R_2 + 2R_1$. This does not change the determinant:
    \[ \begin{vmatrix} a & b & c \\ d & e & f \\ 4a & 4b & 4c \end{vmatrix} \]
    Notice that Row 3 is exactly 4 times Row 1. By the Proportionality Cheat, if one row is a scalar multiple of another, the determinant is $0$. (Alternatively, factor out 4 from Row 3 to get two identical rows, which yields a determinant of 0).
    \textbf{Final Answer:} $0$

    % Solution 11
    \item \textbf{Solution:} A matrix is invertible if and only if its determinant is non-zero. We compute $\det(G)$:
    \[ \det(G) = \begin{vmatrix} 1 & -1 & 2 \\ 2 & k & 4 \\ 3 & -3 & k^2 \end{vmatrix} \]
    Apply $R_2 \leftarrow R_2 - 2R_1$ and $R_3 \leftarrow R_3 - 3R_1$:
    \[ = \begin{vmatrix} 1 & -1 & 2 \\ 0 & k+2 & 0 \\ 0 & 0 & k^2-6 \end{vmatrix} \]
    The matrix is upper triangular, so the determinant is the product of the diagonal: $1(k+2)(k^2-6)$.
    For the matrix to be invertible, $\det(G) \neq 0$.
    Thus, $k+2 \neq 0 \implies k \neq -2$, and $k^2-6 \neq 0 \implies k \neq \pm\sqrt{6}$.
    \textbf{Final Answer:} $k \neq -2$ and $k \neq \pm\sqrt{6}$

    % Solution 12
    \item \textbf{Solution:} Observe the rows of matrix $H$. Row 1 is $[2, 4, 2]$ and Row 2 is $[1, 2, 1]$. Row 1 is exactly $2$ times Row 2. Because two rows are proportional, the determinant is $0$. By the Invertible Matrix Theorem, since $\det(H) = 0$, the matrix is not invertible.
    \textbf{Final Answer:} Not invertible.

    % Solution 13
    \item \textbf{Solution:} Vectors are linearly independent if the matrix formed by them as columns has a non-zero determinant.
    \[ A = \begin{bmatrix} 1 & 2 & 0 \\ -1 & 1 & -3 \\ 2 & 3 & 1 \end{bmatrix} \]
    Expand along Row 1:
    \[ \det(A) = 1 \begin{vmatrix} 1 & -3 \\ 3 & 1 \end{vmatrix} - 2 \begin{vmatrix} -1 & -3 \\ 2 & 1 \end{vmatrix} + 0 \]
    \[ = 1(1 - (-9)) - 2(-1 - (-6)) = 1(10) - 2(5) = 10 - 10 = 0 \]
    Since the determinant is $0$, the matrix is singular, meaning its columns are linearly dependent.
    \textbf{Final Answer:} Linearly dependent.

    % Solution 14
    \item \textbf{Solution:} Form a matrix with the vectors as columns and set the determinant to $0$ (condition for linear dependence).
    \[ \begin{vmatrix} 1 & 2 & 3 \\ 2 & 3 & 4 \\ 3 & 4 & h \end{vmatrix} = 0 \]
    Apply $R_2 \leftarrow R_2 - 2R_1$ and $R_3 \leftarrow R_3 - 3R_1$:
    \[ \begin{vmatrix} 1 & 2 & 3 \\ 0 & -1 & -2 \\ 0 & -2 & h-9 \end{vmatrix} = 0 \]
    Apply $R_3 \leftarrow R_3 - 2R_2$:
    \[ \begin{vmatrix} 1 & 2 & 3 \\ 0 & -1 & -2 \\ 0 & 0 & h-9 - (-4) \end{vmatrix} = \begin{vmatrix} 1 & 2 & 3 \\ 0 & -1 & -2 \\ 0 & 0 & h-5 \end{vmatrix} = 0 \]
    The determinant is $1(-1)(h-5) = 5-h$. For linear dependence, $5-h = 0 \implies h = 5$.
    \textbf{Final Answer:} $h = 5$

    % Solution 15
    \item \textbf{Solution:} By the multiplicative property of determinants, $\det(AB) = \det(A)\det(B)$.
    $\det(AB) = (4)(-3) = -12$.
    \textbf{Final Answer:} $-12$

    % Solution 16
    \item \textbf{Solution:} When multiplying an $n \times n$ matrix by a scalar $c$, the determinant is multiplied by $c^n$. Here, $A$ is $3 \times 3$ and $c = 2$.
    $\det(2A) = 2^3 \det(A) = 8(5) = 40$.
    \textbf{Final Answer:} $40$

    % Solution 17
    \item \textbf{Solution:} By the inverse property, $\det(A^{-1}) = \frac{1}{\det(A)}$.
    $\det(A^{-1}) = \frac{1}{-2} = -\frac{1}{2}$.
    \textbf{Final Answer:} $-\frac{1}{2}$

    % Solution 18
    \item \textbf{Solution:} Apply multiple properties: scalar multiplication, transpose, and inverse.
    $\det(2 A^T B^{-1}) = 2^3 \det(A^T B^{-1})$ (since matrices are $3 \times 3$)
    $= 8 \det(A^T) \det(B^{-1})$
    $= 8 \det(A) \frac{1}{\det(B)}$
    Substitute the given values: $= 8(3)\left(\frac{1}{6}\right) = 24 / 6 = 4$.
    \textbf{Final Answer:} $4$

    % Solution 19
    \item \textbf{Solution:} $-A$ is the same as scalar multiplication by $-1$. Since $A$ is $5 \times 5$, $\det(-A) = (-1)^5 \det(A) = -1(7) = -7$.
    \textbf{Final Answer:} $-7$

    % Solution 20
    \item \textbf{Solution:} By the multiplicative property, $\det(A^4) = (\det(A))^4$.
    We are given $(\det(A))^4 = 16$. Since $A$ has real entries, its determinant must be a real number.
    Taking the fourth root of both sides yields $\det(A) = 2$ or $\det(A) = -2$.
    \textbf{Final Answer:} $2$ or $-2$

    % Solution 21
    \item \textbf{Solution:} Apply the multiplicative property $\det(XY) = \det(X)\det(Y)$ repeatedly:
    $\det(P^{-1}AP) = \det(P^{-1})\det(A)\det(P)$.
    Since determinants are scalar values, multiplication is commutative. We can rearrange:
    $= \det(P^{-1})\det(P)\det(A)$.
    Using the inverse property, $\det(P^{-1}) = \frac{1}{\det(P)}$.
    $= \left(\frac{1}{\det(P)}\right) \det(P) \det(A) = 1 \cdot \det(A) = \det(A)$.
    \textbf{Final Answer:} The proof is complete.

    % Solution 22
    \item \textbf{Solution:} Take the determinant of both sides of the equation $Q^T Q = I$:
    $\det(Q^T Q) = \det(I)$.
    Using the multiplicative property and the fact that $\det(I) = 1$:
    $\det(Q^T)\det(Q) = 1$.
    Using the transpose property $\det(Q^T) = \det(Q)$:
    $\det(Q)\det(Q) = 1 \implies (\det(Q))^2 = 1$.
    Taking the square root gives $\det(Q) = 1$ or $\det(Q) = -1$.
    \textbf{Final Answer:} $\pm 1$

    % Solution 23
    \item \textbf{Solution:} The matrix $M$ is a block upper triangular matrix of the form $\begin{bmatrix} A & B \\ 0 & C \end{bmatrix}$, where $A = \begin{bmatrix} 3 & 2 \\ 0 & 4 \end{bmatrix}$ and $C = \begin{bmatrix} 2 & 5 \\ 1 & 3 \end{bmatrix}$.
    The determinant of a block upper triangular matrix is the product of the determinants of its diagonal blocks: $\det(M) = \det(A)\det(C)$.
    $\det(A) = (3)(4) - (0)(2) = 12$.
    $\det(C) = (2)(3) - (5)(1) = 6 - 5 = 1$.
    $\det(M) = 12 \times 1 = 12$.
    \textbf{Final Answer:} $12$

    % Solution 24
    \item \textbf{Solution:} The matrix $N$ is a block lower triangular matrix of the form $\begin{bmatrix} A & 0 \\ B & C \end{bmatrix}$, where $A = [3]$ (a $1 \times 1$ matrix) and $C = \begin{bmatrix} 1 & 2 \\ 3 & 4 \end{bmatrix}$.
    $\det(N) = \det(A)\det(C)$.
    $\det(A) = 3$.
    $\det(C) = (1)(4) - (2)(3) = 4 - 6 = -2$.
    $\det(N) = 3 \times (-2) = -6$.
    \textbf{Final Answer:} $-6$

    % Solution 25
    \item \textbf{Solution:} This is a $3 \times 3$ Vandermonde matrix with $x_1 = 2$, $x_2 = 3$, and $x_3 = 5$.
    The formula for the Vandermonde determinant is $\det(V) = (x_2 - x_1)(x_3 - x_1)(x_3 - x_2)$.
    $\det(V) = (3 - 2)(5 - 2)(5 - 3) = (1)(3)(2) = 6$.
    \textbf{Final Answer:} $6$

    % Solution 26
    \item \textbf{Solution:} This is a Vandermonde matrix where $x_1 = 1$, $x_2 = 1$, and $x_3 = 4$.
    Because $x_1 = x_2$, the first two columns are identical. By the proportionality cheat (or the Vandermonde formula where $(x_2 - x_1) = 0$), the determinant is exactly $0$.
    \textbf{Final Answer:} The determinant is $0$ because two columns are identical ($x_1 = x_2$).

    % Solution 27
    \item \textbf{Solution:} \textbf{Flaw:} The determinant is not a linear operator over matrix addition; $\det(A+B) \neq \det(A) + \det(B)$. This is the Linearity Trap.
    \textbf{Correct Solution:} First, compute the matrix sum $A+B = \begin{bmatrix} 2+3 & 0 \\ 0 & 2+3 \end{bmatrix} = \begin{bmatrix} 5 & 0 \\ 0 & 5 \end{bmatrix}$.
    The determinant of this diagonal matrix is $5 \times 5 = 25$.
    \textbf{Final Answer:} $25$

    % Solution 28
    \item \textbf{Solution:} We are given $A^T = -A$. Take the determinant of both sides:
    $\det(A^T) = \det(-A)$.
    By the transpose property, $\det(A^T) = \det(A)$.
    By the scalar multiplication property for an $n \times n$ matrix, $\det(-A) = (-1)^n \det(A)$.
    Equating the two gives: $\det(A) = (-1)^n \det(A)$.
    Since $n$ is odd, $(-1)^n = -1$.
    Therefore, $\det(A) = -\det(A)$.
    Adding $\det(A)$ to both sides yields $2\det(A) = 0$, which implies $\det(A) = 0$.
    \textbf{Final Answer:} The proof is complete.

    % Solution 29
    \item \textbf{Solution:} We are given $A^k = 0$. Take the determinant of both sides:
    $\det(A^k) = \det(0)$.
    The determinant of the zero matrix is $0$.
    Using the multiplicative property repeatedly, $\det(A^k) = (\det(A))^k$.
    So, $(\det(A))^k = 0$.
    The only real (or complex) number that yields $0$ when raised to a positive integer power is $0$. Therefore, $\det(A) = 0$.
    \textbf{Final Answer:} The proof is complete.

    % Solution 30
    \item \textbf{Solution:} \textbf{Flaw:} The operation $R_2 \leftarrow 3R_1 - R_2$ is not a standard row replacement. It is a combination of scaling Row 2 by $-1$ and then adding $3R_1$. Scaling a row by $-1$ multiplies the determinant by $-1$. The flawed solution ignored this scaling effect.
    \textbf{Correct Solution:} The standard replacement that preserves the determinant is $R_2 \leftarrow R_2 - 3R_1$.
    Applying $R_2 \leftarrow R_2 - 3R_1$ to $A = \begin{bmatrix} 1 & 2 \\ 3 & 4 \end{bmatrix}$ gives $\begin{bmatrix} 1 & 2 \\ 0 & -2 \end{bmatrix}$.
    The determinant is $1 \times (-2) = -2$. (Check: $\det(A) = (1)(4) - (2)(3) = 4 - 6 = -2$).
    \textbf{Final Answer:} $-2$

\end{enumerate}

\newpage

\section{Concept Checklist \& Question Mapping (Appendix)}
\begin{itemize}
    \item \textbf{C1: Determinant via Row Reduction (Type 1)} - Questions 1, 2, 3
    \item \textbf{C2: Hybrid Evaluation (Type 5)} - Questions 4, 5, 6
    \item \textbf{C3: Abstract Row Operations (Type 2)} - Questions 7, 8, 9, 10
    \item \textbf{C4: Invertibility via Determinants (Type 3)} - Questions 11, 12
    \item \textbf{C5: Linear Independence via Determinants (Type 3)} - Questions 13, 14
    \item \textbf{C6: Multiplicative Property (Type 4)} - Questions 15, 21
    \item \textbf{C7: Transpose Property (Type 4)} - Question 18
    \item \textbf{C8: Inverse Property (Type 4)} - Questions 17, 18
    \item \textbf{C9: Scalar Multiplication of a Matrix (Type 4)} - Questions 16, 18, 19
    \item \textbf{C10: Determinants of Block Matrices (Type 6)} - Questions 23, 24
    \item \textbf{C11: The Vandermonde Determinant (Type 7)} - Questions 25, 26
    \item \textbf{C12: Theoretical Proofs \& Flaw Finding (Type 8)} - Questions 20, 22, 27, 28, 29, 30
\end{itemize}

\end{document}